\chapter{Bounded Completeness: Multi-Hazard Corner Cases (Phase 2.5)}
\label{chap:p2-5}

\section{Goal}

The proposal's most distinctive claim was that a good explanation should capture \emph{every} hazard an
image exhibits, not merely one. The pilot could not test this: priority-collapse labeling (fixed in
Phase~1.2) assigned each image a single class, so an image that violated two rules was only ever tested
against one. The plan (item 4.5) asks for three things here: bake the worker-loss correction into the
headline completeness number (already delivered in Phase~1.4); state explicitly that completeness is
not an independent confirmation of descriptive accuracy; and --- the substantive new work --- test
multi-hazard completeness on images that violate two or more rules, reporting it as its own result.

\section{What Was Done}

\subsection{Worker-loss correction and the shared-construction caveat}

The completeness headline is already worker-loss-correctable: Phase~1.4 added
\texttt{classify\_sample\_worker\_loss\_corrected} and reduced the pilot's ``supported'' rate from
43.6\% to 41.1\% as a pure rollup, no reruns. Phase~2.5 records the caveat the plan asks for plainly:
the bounded-completeness verdict is computed \emph{from} the descriptive-accuracy masking results (it is
a rollup of the same top-1/top-2 answer-flip signal), so a sample being both ``descriptively accurate''
and ``complete'' is one measurement reported twice, not two independent confirmations. The two metrics
should be cited together with that dependence stated, never as mutually corroborating.

\subsection{Multi-hazard completeness}

A new metric, \texttt{multi\_hazard\_verdict}, classifies an image as \texttt{complete} (every violated
hazard's own region is load-bearing), \texttt{partial}, or \texttt{none}. A focused script
(\texttt{12\_multihazard\_completeness.py}) finds every test-split image violating two or more rules and,
for each violated rule \emph{independently}, masks that rule's rule-aware top region (Phase~1.1) and
checks whether the rule's answer flips, worker-loss-corrected (Phase~1.4). This is the first time the
pipeline tests one image against several rules --- the \texttt{(image, rule)} multi-testing the full
Phase~5 scale run generalizes, scoped here to the multi-hazard subset.

\section{Results}

The test split contains 22 images that violate two or more rules (most commonly rule\_1+rule\_4, nine
images, and rule\_1+rule\_3, five). Testing each against all of its hazards:

\begin{center}
\begin{tabular}{p{4.6cm}p{2.4cm}p{2.4cm}}
\toprule
\textbf{verdict} & \textbf{count} & \textbf{share} \\
\midrule
complete (all hazards' regions load-bearing) & 6 & 27.3\% \\
partial (some, not all)                      & 10 & 45.5\% \\
none (no hazard's region load-bearing)       & 6 & 27.3\% \\
\bottomrule
\end{tabular}
\captionof{table}{Multi-hazard completeness over the 22 test-split images that violate two or more
rules. Only about a quarter have an explanation that captures every hazard; on average 1.05 of the 2.09
hazards per image are supported.}
\label{tab:p25-multihazard}
\end{center}

The result is a direct, and sobering, measurement of the proposal's core promise: for multi-hazard
scenes, Florence-2's grounding-proxy explanation captures only about \emph{half} the hazards present
(1.05 of 2.09 on average), and fewer than one image in three has an explanation that is complete across
all its hazards. The proxy reliably points at one hazard --- typically the most visually salient object
--- and leaves the others unexplained. This is exactly the failure mode a single-label evaluation could
never surface, because it never asked the image about its second hazard. It is also a clean motivation
for the native-VQA models of Phase~5: a model that reasons about the whole scene, rather than grounding
one phrase at a time, is the natural candidate to raise multi-hazard completeness.

\subsection{Higher-order necessity}

The plan flags an optional extension: a region necessary only \emph{in combination} with another. The
pilot's cumulative masking (top-1, then top-1+top-2) plus the Phase~2.1 \texttt{non\_monotonic} flag
already detect the simplest form of this --- a region whose removal only matters alongside another's ---
so it is partially in place. A full higher-order necessity analysis (every subset of regions) is a
combinatorial extension deferred to the scale run, where the multi-worker, multi-object detection of D4
makes more than two candidate regions per sample common enough to be worth the extra masking.

\section{An Honest Note on Sample Size}

Twenty-two multi-hazard images is a pilot-scale floor, not a powered estimate; the 27.3\% complete rate
should be read as an existence result (multi-hazard incompleteness is real and common) rather than a
precise rate. The plan anticipates exactly this and provides the remedy: the train split supplies
additional multi-hazard examples, and the scale run stratifies and oversamples them to reach an adequate
$n$ per rule-combination. The machinery to test them --- the metric, the per-rule masking, the
worker-loss correction --- is in place and validated here; only the sample count grows.

\section{Standard vs. Non-Standard Choices}

\textbf{Standard.} Testing each hazard's region for necessity independently, and rolling the per-hazard
verdicts into a per-image completeness class, is the natural construction once multi-label data exists.

\textbf{Non-standard, and the contribution.} The phase reports a metric the pilot's own design made
unmeasurable, and reports it \emph{against} the project's interest --- 27.3\% complete is a weak result
for the proxy, stated plainly rather than buried. It also states the completeness/descriptive-accuracy
dependence explicitly, refusing to present a rollup of one signal as two independent confirmations ---
a common way evaluation sections overstate their evidence.

\section{Implications}

Multi-hazard completeness is now a measured quantity, and its low value (27.3\% complete) is one of the
strongest arguments in the study for moving beyond the single-phrase grounding proxy to a scene-level
native-VQA model (Phase~5). The per-rule masking harness built here is the same one the scale run uses
to test every image against every rule it violates, so this phase doubles as the working prototype of
that consumption path. The immediate carry-forward is sample size: the test is ready; the multi-hazard
$n$ needs to grow from the train split before the rate is quoted as more than an existence result.
